\section{Object-Relational Mapping}

% ----------------------------------------------------------
\begin{frame}{Object-Relational Mapping (ORM) -- Introduction}
  \begin{block}{Definition}
    An \textbf{ORM (Object-Relational Mapper)} is a software library that automatically translates
    between the relational database world (tables, rows, columns) and the object-oriented programming
    world (classes, instances, fields), allowing developers to work with database data using their
    programming language's idioms rather than writing raw SQL.
  \end{block}

  \begin{center}
  \renewcommand{\arraystretch}{1.25}
  \begin{tabular}{@{}ll@{}}
    \toprule
    \textbf{Relational Database} & \textbf{Object-Oriented (Java)} \\
    \midrule
    Relation (Table)      & Class \\
    Tuple (Row)           & Instance (object) of the class \\
    Attribute (Column)    & Field / property \\
    Primary Key           & \texttt{id} field (auto-generated) \\
    Foreign Key           & Object reference to another class instance \\
    JOIN query            & Navigating an object reference (\texttt{obj.getX()}) \\
    Stored Procedure      & Method on repository or service class \\
    \bottomrule
  \end{tabular}
  \end{center}

  \begin{exampleblock}{Popular ORM Frameworks}
    \textbf{Hibernate / JPA} (Java) \quad \textbf{SQLAlchemy} (Python) \quad
    \textbf{ActiveRecord} (Ruby) \quad \textbf{Eloquent} (PHP) \quad
    \textbf{Entity Framework} (.NET)
  \end{exampleblock}
\end{frame}

% ----------------------------------------------------------
\begin{frame}[fragile]{The Object-Relational Impedance Mismatch}
  Relational databases and object-oriented languages represent data very differently.
  This fundamental tension is called the \textbf{object-relational impedance mismatch}.

  \begin{columns}[T]
    \column{0.48\linewidth}
      \textbf{Relational Model}
      \begin{itemize}
        \item Data in flat tables with rows
        \item Relationships via foreign keys / JOINs
        \item No identity concept (rows have no memory address)
        \item Set-based operations
        \item No inheritance
      \end{itemize}

    \column{0.48\linewidth}
      \textbf{Object-Oriented Model}
      \begin{itemize}
        \item Data in nested object graphs
        \item Relationships via object references
        \item Identity via heap address / \texttt{==}
        \item Object-at-a-time operations
        \item Inheritance, polymorphism
      \end{itemize}
  \end{columns}

  \vspace{6pt}
  \begin{alertblock}{The Mismatch in Practice}
    An \texttt{Artist} object may contain a \texttt{List<Album>}, each \texttt{Album} a
    \texttt{List<Track>}. In the DB these are three separate tables. The ORM must bridge this
    gap automatically -- and this is where most ORM complexity (and bugs) originates.
  \end{alertblock}
\end{frame}

% ----------------------------------------------------------
\begin{frame}{[GRAPHIC PLACEHOLDER] Impedance Mismatch Diagram}
  \begin{alertblock}{INSERT GRAPHIC HERE}
    Add a diagram showing: left side = Java object graph (Artist \textrightarrow{} List<Album>
    \textrightarrow{} List<Track>); right side = three relational tables with FK arrows;
    middle = ORM bridging the two worlds.
  \end{alertblock}
\end{frame}
